\section{Pointwise smoothing and conditioned grid entropy}
\label{sec:smoothed-grid}\label{rec:sec-smoothing}

We now construct line densities whose smoothing estimate holds at each
permutation. It can therefore be used after the entire row--column grid
has been conditioned on compatibility, provided each auxiliary subset
retains its Bernoulli law independently of the completed array. This pointwise estimate
is the additional output of the present route. Its final indexed
singular-value bound has the same form as that of
Section~\ref{sec:band-grid}.

For the final recursion, put $n=2^d$,
$a=2^{\lfloor d/2\rfloor}$ and $b=2^{\lceil d/2\rceil}$.
Write $T_n=YX$, where $X$ is the tensor product of row sweeps
on the $a$ by $b$ grid and $Y$ the tensor product of column sweeps.
The pointwise construction first treats one line. We then use it in
a complete-grid entropy exposure, prove a separate sparse estimate
by chronological mergers, and combine the two estimates at a split.

\subsection{A pointwise density construction}

\begin{lemma}\label{rec:pointwise-smoothing}
Let $\mu\vdash q$, $j=q-\mu_1$, and $\ell=\log m$, where
$m\le q\le2m$ and $m$ is sufficiently large. Let $\mathcal C_\mu$
be the space spanned by the matrix coefficients of $V_\mu$, with
the $L^2$ inner product for uniform measure on $S_q$. Put
\[
 p_t=e^{-\sqrt\ell-t},\quad
 0\le t\le\lfloor\ell/32-\sqrt\ell\rfloor,
 \qquad h_0=\lfloor\ell^{-4}j\rfloor.
\]
For each $a\in\mathcal C_\mu$ with $\mathbb E|a|^2=1$, there is a
positive sum $Y$ of squared functions in $\mathcal C_\mu$ such that
$\mathbb E Y=1$, $|a(g)|^2\le2Y(g)$, and
$Y(g)\le D_\mu^2$ for every $g\in S_q$. For every allowed $t$,
every $g$ with $Y(g)>0$, and a Bernoulli-$p_t$ subset $A$ of sites,
\begin{equation}\label{rec:smoothing-ratio}
 \E_A\log\frac{Y(g)}{(Q_A Y)(g)}
 \le\log4+C\ell^{-4}j\log q.
\end{equation}
Here $(Q_A Y)(g)=|S_A|^{-1}\sum_{h\in S_A}Y(gh)$, where $S_A$
fixes every site outside $A$. The same assertion holds if the
initial $|a|^2$ is any positive sum of squared coefficient functions
of mean one.
\end{lemma}
\begin{proof}
Let $Z_A$ be the projection onto restriction types of level greater
than $h_0$.  Averaging over $A$ commutes with the full symmetric group,
so $\E Z_A=c_tI$.  We identify the finite comparison underlying the
bound for $c_t$.  Let $\mathcal S$ be the increasing lists
$(a_1<\cdots<a_j)$ in $\{1,\ldots,q\}$, ordered coordinatewise.
Coordinatewise minimum and maximum are again increasing lists;
thus $\mathcal S$ is a finite distributive lattice.  Its uniform
weight is log-supermodular (the lattice inequality is equality).
The finite FKG association inequality therefore applies to this
uniform measure \cite[Proposition~1]{FortuinKasteleynGinibre1971}.

Fix a relative standard ordering of the $j$ boxes below the first
row.  For a chosen list in $\mathcal S$, put its entries in those
boxes in the fixed order and put the complementary entries increasingly
in the first row.  The event that this filling is a standard tableau
is increasing on $\mathcal S$: moving bottom entries later preserves
all comparisons within the lower diagram, moves each complementary
first-row entry earlier, and hence preserves every comparison from
the first row to the second row.  Conditional on this relative lower
ordering, a uniform standard tableau is exactly a uniform list
conditioned on that increasing event.  The event that at least $h$
bottom entries occur among $1,\ldots,b$ is decreasing on $\mathcal S$.
FKG, applied to its complement and the increasing validity event,
shows that conditioning on validity can only lower this early-count
tail.  Averaging over the possible relative lower orderings preserves
the inequality.

By branching, conditional on $|A|=b$, the restriction level equals
the number of bottom entries among the first $b$ entries of that
uniform tableau.  Its upper tail is therefore bounded by the count
for a uniform $j$-subset of $\{1,\ldots,q\}$.  Averaging
$b\sim\operatorname{Binomial}(q,p_t)$ turns the latter count into
$\operatorname{Binomial}(j,p_t)$.  With $k_0=h_0+1$, its tail is at
most $(ejp_t/k_0)^{k_0}$ when $k_0\le j$, and is zero otherwise.
Since $jp_t/k_0\le\ell^4e^{-\sqrt\ell}$, this is at most
$\ell^{-10}$ for all sufficiently large $\ell$, uniformly in $j,t$.
This proves the asserted bound for $c_t$.

We next turn this small average high-level projection into an
estimate valid at each permutation. Right translations act on
$\mathcal C_\mu$ as copies of $V_\mu$; use the corresponding
projection $Z_A$ on this coefficient space. The identity
$\mathbb E_A Z_A=c_tI$ is still valid, including all copies. For a
positive operator $W$ on $\mathcal C_\mu$, define
\[
 \mathcal A_t(W)=\mathbb E_A Z_AWZ_A,
 \qquad \mathcal A(W)=\ell^4\sum_t\mathcal A_t(W).
\]
There are $O(\ell)$ values of $t$, and $c_t\le\ell^{-10}$.
Consequently $\operatorname{Tr}\mathcal A(W)\le\frac12
\operatorname{Tr}W$ for sufficiently large $\ell$. The positive
series
\[
 W'=\sum_{r\ge0}\mathcal A^r(W)
\]
converges in trace norm, satisfies $W'\ge W$, has trace at most
$2\operatorname{Tr}W$, and obeys
$\ell^4\mathcal A_t(W')\le W'$ for every $t$.

Here is the concrete interpretation of these positive operators.
For $g\in S_q$, let $e_g\in\mathcal C_\mu$ represent evaluation
at $g$. Orthogonality gives
$\mathbb E_g e_ge_g^*=I$ and $\|e_g\|^2=\dim\mathcal C_\mu
=D_\mu^2$. The function $f_W(g)=\langle e_g,We_g\rangle$
is a positive sum of squared coefficient functions, and
$\mathbb E f_W=\operatorname{Tr}W$. Start with the rank-one
operator $W=aa^*$, or with the positive operator representing the
given sum, so that $\operatorname{Tr}W=1$. Set
\[
 \widetilde W=W'/\operatorname{Tr}W',\qquad
 Y(g)=\langle e_g,\widetilde W e_g\rangle.
\]
The trace bound proves the majorization and normalization in the
statement, and $Y(g)\le\|e_g\|^2$ proves its pointwise upper bound.

Fix $g$ with $Y(g)>0$. Write $Y_{\rm hi,A}$ and $Y_{\rm lo,A}$
for the positive functions obtained from $\widetilde W$ by compression
with $Z_A$ and $I-Z_A$. The operator inequality above gives
\[
 \mathbb E_A Y_{\rm hi,A}(g)\le\ell^{-4}Y(g).
\]
Thus $Y_{\rm hi,A}(g)\le Y(g)/4$ except on an event of probability
at most $4\ell^{-4}$. The scalar inequality $|u+v|^2\le
2|u|^2+2|v|^2$, applied to a spectral decomposition of
$\widetilde W$, gives $Y\le2Y_{\rm lo,A}+2Y_{\rm hi,A}$.
On the complementary event, $Y_{\rm lo,A}(g)\ge Y(g)/4$.

We spell out the evaluation estimate used here. On $S_A$, every
coefficient of restriction level at most $h$ lies in the sum of the
coefficient spaces of representations occurring on ordered
$\min(h,|A|)$-tuples. That sum has dimension at most $q^{2h}$.
For a translation-invariant coefficient space of dimension $D$,
Cauchy--Schwarz at the identity bounds evaluation squared by $D$
times its uniform squared norm. Applying this fact to each squared
coefficient in $Y_{\rm lo,A}$ gives
\[
 (Q_A Y)(g)\ge (Q_A Y_{\rm lo,A})(g)
 \ge q^{-2h_0}Y_{\rm lo,A}(g).
\]
The first inequality uses orthogonality of the high- and low-level
parts after averaging over $S_A$; their cross terms vanish. Every
restriction of $V_\mu$ has level at most $j$, so the same evaluation
argument without truncation always gives
$(Q_A Y)(g)\ge q^{-2j}Y(g)$. Therefore
\[
 \mathbb E_A\log\frac{Y(g)}{(Q_A Y)(g)}
 \le\log4+2h_0\log q+8\ell^{-4}j\log q,
\]
which is \eqref{rec:smoothing-ratio}.

Every estimate in this last paragraph holds for each fixed $g$.
In particular one may subsequently condition the grid permutation
on compatibility, or average $g$ under any other law, provided the
auxiliary subset $A$ retains its independent Bernoulli law. The
bound is uniform in that changed law of $g$. This is the
conditioning invariant supplied by the positive series.
\end{proof}

\subsection{Entropy after conditioning on a complete grid}

The pointwise estimate now controls the line-density terms in an
entropy calculation under the law conditioned on compatibility.
The cell priorities used to expose that law will remain independent
of the completed array.

\begin{lemma}\label{rec:smoothed-grid-fourth}
Let $m=\min(a,b)\ge2$, with $a,b\le2m$. For product Fourier projections $I,J$ in the line groups, of local
types $\mu_i$ with at most half as many rows as line sites and carrier
ranks $r_i$, a zero local rank makes $I$ or $J$ zero. Otherwise put
$D_0=\sum_i\log(D_{\mu_i}r_i)$. Then
\[
 \Tr_{\mathrm{reg}}|JI|^4
 \le\exp\{(1+C/\sqrt{\log m})D_0+Cn m^{-1/40}\}.
\]
\end{lemma}
\begin{proof}
Write $K=(S_b)^a$ for the row group and $L=(S_a)^b$ for the
column group, where $n=ab$ and $m\le a,b\le2m$. The trace is the
unnormalized trace on the regular representation of $S_n$. We prove
the estimate for sufficiently large $m$; the finitely many sizes
$2\le m<m_0$ are covered by increasing the absolute constant $C$.
Let $\mathcal L$ be the set of all $a+b$ lines. For
$\gamma\in\mathcal L$, denote its length by $q_\gamma$, its
Fourier type by $\mu_\gamma$, its dimension by
$d_\gamma=D_{\mu_\gamma}$, and its carrier projection by
$P_\gamma$. Its rank is $r_\gamma$. A zero rank makes $I$ or $J$
zero, so we assume $1\le r_\gamma\le d_\gamma$. The hypotheses
are $m\le q_\gamma\le2m$ and
$\ell(\mu_\gamma)\le q_\gamma/2$. Set
\[
 j_\gamma=q_\gamma-(\mu_\gamma)_1,
 \qquad \ell=\log m,
 \qquad D_0=\sum_{\gamma\in\mathcal L}
                         \log(d_\gamma r_\gamma).
\]
Here $\ell(\mu_\gamma)$ denotes the number of rows of the partition;
the unadorned symbol $\ell$ denotes $\log m$.

We first identify the probability whose entropy will be estimated.
Choose unitary realizations $\rho_\gamma$ and put
\[
 f_\gamma(g)=d_\gamma
       \operatorname{Tr}\bigl(P_\gamma\rho_\gamma(g^{-1})\bigr),
 \qquad
 \eta_\gamma(g)=\frac{|f_\gamma(g)|^2}
                            {d_\gamma r_\gamma}.
\]
All averages on a line group are with respect to its uniform
probability measure $U_\gamma$. Matrix-coefficient orthogonality
and $P_\gamma=P_\gamma^*=P_\gamma^2$ give
\begin{equation}\label{rec:smoothed-line-normalization}
 \mathbb E_{U_\gamma}|f_\gamma|^2=d_\gamma r_\gamma,
 \qquad
 f_\gamma(g^{-1})=\overline{f_\gamma(g)},
 \qquad \mathbb E_{U_\gamma}\eta_\gamma=1.
\end{equation}
Thus $\eta_\gamma$ is an inversion-invariant probability density.
For each local projection, the regular rank and squared coefficient
mass in Lemma~\ref{lem:coefficient-compatibility} are both
$d_\gamma r_\gamma$. Applying that lemma to $I,J$ gives
\begin{equation}\label{rec:smoothed-trace-probability}
 \operatorname{Tr}_{\rm reg}|JI|^4
 \le e^{D_0}\frac{n!}{|K||L|}\,\mathbb P_\eta(\mathcal S).
\end{equation}
We now describe $\mathcal S$ and $\mathbb P_\eta$ explicitly.

Let $r=(r_i)\in K$ and $c=(c_t)\in L$ have the independent
product line law $\mathbb P_\eta$. At the intermediate cell
$x=(i,t)\in[a]\times[b]$, write
\[
 X_{it}=r_i^{-1}(t),\qquad Z_{it}=c_t(i).
\]
The vector $X_i=(X_{it})_{t\in[b]}$ is a permutation of $[b]$,
and $Z_t=(Z_{it})_{i\in[a]}$ is a permutation of $[a]$.
The event $\mathcal S$ is that the $n$ pairs
$(X_{it},Z_{it})$ are all different. Equivalently they give a
bijection from the intermediate cells to $[b]\times[a]$.
Indeed the row-then-column permutation $cr$ admits a
column-then-row factorization exactly when, for each initial column
and each final row, there is exactly one such intermediate cell.
This is the stated distinctness condition. Inversion invariance in
\eqref{rec:smoothed-line-normalization} ensures that the use of
$r_i^{-1}$ in the row variable preserves its line density.
Under $\mathbb P_\eta$ all the line permutations $X_i,Z_t$
are independent, with densities $\eta_\gamma$ relative to
$U_\gamma$. This also specifies the orientation of every line
variable used below.

Apply Lemma~\ref{rec:pointwise-smoothing} to
$f_\gamma/(d_\gamma r_\gamma)^{1/2}$. It gives a density
$Y_\gamma$ satisfying
\begin{equation}\label{rec:smoothed-line-densities}
 \eta_\gamma\le2Y_\gamma,
 \qquad \mathbb E_{U_\gamma}Y_\gamma=1,
 \qquad Y_\gamma\le d_\gamma^2,
\end{equation}
and the pointwise estimate \eqref{rec:smoothing-ratio} with
$j=j_\gamma$ and $q=q_\gamma$. For the trivial type
$\mu_\gamma=(q_\gamma)$ take $Y_\gamma=\eta_\gamma=1$,
so no factor two is needed on that line. Let $N_*$ be the number
of nontrivial line types and define the product law
\[
 \mathbb M=\bigotimes_{\gamma\in\mathcal L}
                                  (Y_\gamma U_\gamma)
 \quad\hbox{on}\quad
 \Omega=\prod_{i=1}^a S_b\times\prod_{t=1}^b S_a.
\]
Then
\begin{equation}\label{rec:smoothed-majorization-cost}
 \mathbb P_\eta(\mathcal S)\le2^{N_*}\mathbb M(\mathcal S).
\end{equation}
Put $z=\mathbb M(\mathcal S)$. If $z=0$ the assertion follows
already. Otherwise set $\mathbb N=\mathbb M(\,\cdot\mid\mathcal S)$.
The remaining task is to show that this conditioning costs almost
$n$ units of relative entropy, after a small charge for the line
densities.

Independently of the array, assign each cell $x$ a uniform priority
$\tau_x\in(0,1)$ and reveal the pair $(X_x,Z_x)$ in increasing
priority order. Ties have probability zero. Conditional on the
complete vector of priorities and the values already revealed,
let $q_x$ be the prediction of the next pair under $\mathbb N$,
and let $\pi_x$ be its prediction under $\mathbb M$. If $x=(i,t)$,
let $B_x$ be the symbols of $[b]$ not yet revealed in row $i$ and
$C_x$ the symbols of $[a]$ not yet revealed in column $t$.
Write $R_x=|B_x|$ and $S_x=|C_x|$. The uniform product law on
$\Omega$ predicts the reference distribution
\[
 \pi_{0,x}(u,v)=\frac1{R_xS_x},\qquad (u,v)\in B_x\times C_x.
\]
The product structure of $\mathbb M$ implies that $\pi_x$ is the
product of the two line predictions, one for row $i$ and one for
column $t$. In particular, it depends only on the revealed values
in those two lines, even though the conditioning specifies the
entire past. Let $A_x\subseteq B_x\times C_x$ consist of the
pairs that have not appeared at any previously revealed cell.
The prediction $q_x$ is supported on $A_x$ and is absolutely
continuous with respect to $\pi_x$. The entropy chain rule gives
the exact equality
\begin{equation}\label{rec:smoothed-exposure-KL}
 -\log z
 =\mathbb E_{\mathbb N,\tau}\sum_x
                              \operatorname{KL}(q_x\Vert\pi_x).
\end{equation}

We use a slowly changing interpolation with the uniform reference
law. Put
\[
 T=\left\lfloor\ell/32-\sqrt\ell\right\rfloor,
 \qquad p_t=e^{-\sqrt\ell-t}\quad(0\le t\le T),
 \qquad p_{\min}=m^{-1/32},
\]
and define, for $0<p<1$,
\begin{equation}\label{rec:smoothed-epsilon-schedule}
 \epsilon(p)=\frac8{\sqrt\ell}
              +\frac8\ell\sum_{t=0}^T\mathbf1_{\{p<p_t\}}.
\end{equation}
For sufficiently large $m$, $0<\epsilon(p)<1/2$. Moreover
\begin{equation}\label{rec:smoothed-epsilon-bound}
 p^{-1/\epsilon(p)}\le m^{1/8}
                 \qquad(p\ge p_{\min}).
\end{equation}
To check this, put $s=-\log p$. For $s\le\sqrt\ell$ the
ratio $s/\epsilon(p)$ is at most $\ell/8$. In each later
interval the number of crossed thresholds is an integer $h$ with
$s\le\sqrt\ell+h$, while
$\epsilon(p)=8(\sqrt\ell+h)/\ell$. This remains true in the
last interval up to $s=\ell/32$, by the definition of $T$.

At cell $x$, put $p=1-\tau_x$ and abbreviate
$\epsilon=\epsilon(p)$. Define the mixed available mass
\[
 Z_x^*=\sum_{v\in A_x}
                   \pi_x(v)^{1-\epsilon}\pi_{0,x}(v)^\epsilon.
\]
It is positive on every history with positive $\mathbb N$
probability. Normalizing the summands on $A_x$ and taking their
relative entropy from $q_x$ yields
\begin{equation}\label{rec:smoothed-mixed-KL}
 \operatorname{KL}(q_x\Vert\pi_x)
 \ge-\log Z_x^*
       -\epsilon(p)\,\mathbb E_{q_x}
                        \log\frac{\pi_x}{\pi_{0,x}}.
\end{equation}
This is an identity with a nonnegative relative-entropy term
discarded; it does not replace $\pi_x$ by a conditioned line
marginal. H\"older's inequality also gives $Z_x^*\le1$.

We next bound the first term in
\eqref{rec:smoothed-mixed-KL} uniformly over the successful array.
Fix a full array $g\in\mathcal S$ with $\mathbb M(g)>0$,
the current priority $\tau_x=1-p$, and all priorities in the
current row and column. The two predictions $\pi_x,\pi_{0,x}$
and their common candidate set $B_x\times C_x$ are now fixed.
Each candidate pair $v$ occurs at a unique cell $y(v)$ in $g$.
Call it special if $y(v)$ is in the current row or column; there
are at most $a+b-1$ such pairs. For every other candidate,
$v\in A_x$ exactly when $\tau_{y(v)}\ge\tau_x$, an event of
conditional probability $p$. These outside priorities remain
independent uniform variables under the present conditioning.
No independence between different blocking events is needed.
Thus, writing $E_x$ for the set of special candidates,
\begin{align*}
 \mathbb E_{\rm outside}Z_x^*
 &\le\sum_{v\in B_x\times C_x}
          \pi_x(v)^{1-\epsilon}\pi_{0,x}(v)^\epsilon c(v),\\
 c(v)&=\begin{cases}1,&v\in E_x,\\p,&v\notin E_x.\end{cases}
\end{align*}
H\"older, with exponents $1/(1-\epsilon)$ and $1/\epsilon$,
therefore gives the precise mixed-mass estimate
\begin{equation}\label{rec:smoothed-competitor-mass}
 \mathbb E_{\rm outside}Z_x^*
 \le\left(\sum_v\pi_{0,x}(v)c(v)^{1/\epsilon}\right)^\epsilon
 \le\bigl(p^{1/\epsilon}+\pi_{0,x}(E_x)\bigr)^\epsilon.
\end{equation}
The first H\"older factor is one because $\sum_v\pi_x(v)=1$.

Conditional on $g$ and $p$ alone, the two available counts have
laws
\[
 R_x=1+\operatorname{Binomial}(b-1,p),\qquad
 S_x=1+\operatorname{Binomial}(a-1,p).
\]
The priorities other than $\tau_x$ in these two lines are
independent. A binomial lower-tail bound shows that
\[
 \mathbb P_\tau\{R_x<bp/2\ \hbox{or}\ S_x<ap/2\mid g,p\}
                                        \le C e^{-cpm}.
\]
On the complementary event, uniformity of $\pi_{0,x}$ gives
$\pi_{0,x}(E_x)\le C/(mp^2)$. For $p\ge p_{\min}$,
\eqref{rec:smoothed-epsilon-bound} consequently implies
\begin{align*}
 \log\mathbb E_{\rm outside}Z_x^*
 &\le\log p+
       \epsilon\log\left(1+\frac{C p^{-1/\epsilon}}{mp^2}\right)\\
 &\le\log p+C m^{-13/16}
 \le\log p+C m^{-1/2}.
\end{align*}
On the exceptional count event we still have $\log Z_x^*\le0$.
Jensen's inequality for the outside priorities, followed by
averaging the priorities on the two lines, now gives
\[
 \mathbb E_\tau[\log Z_x^*\mid g,p]
 \le\log p+C m^{-1/2}+C|\log p|e^{-cpm}
 \le\log p+C m^{-1/2}
                    \qquad(p\ge p_{\min}).
\]
Here $pm\ge m^{31/32}$ absorbs the last term, uniformly in this
range. For $p<p_{\min}$ use only $-\log Z_x^*\ge0$. Since
$p=1-\tau_x$ is uniform independently of $g$, integration gives
\begin{equation}\label{rec:smoothed-slot-saving}
 \mathbb E_\tau[-\log Z_x^*\mid g]
 \ge\int_{p_{\min}}^1(-\log p)\,dp-Cm^{-1/2}
 \ge1-Cm^{-1/40}.
\end{equation}
This bound holds for every such completed successful array and
therefore also after averaging under $\mathbb N$.

It remains to account for the second term in
\eqref{rec:smoothed-mixed-KL}. We do this by an exact telescoping
identity on each line. Fix a line $\gamma$ and its completed
permutation $g_\gamma$. Its domain sites carry the priorities of
the corresponding cells. For $0\le u\le1$, let
\[
 A_\gamma(u)=\{v:\tau_v>u\},
 \qquad
 F_\gamma(u)=
       (Q_{A_\gamma(u)}Y_\gamma)(g_\gamma),
 \qquad H_\gamma(u)=\log F_\gamma(u).
\]
The quantity $F_\gamma(u)$ is exactly the density, relative to
the uniform line law, of the images already exposed on that line.
In fact, permutations agreeing with $g_\gamma$ on the exposed
domain sites form the right coset $g_\gamma S_{A_\gamma(u)}$;
averaging $Y_\gamma$ over its uniform completions is precisely
the displayed $Q_A$ average. This verifies the orientation of
the averaging operator required by
Lemma~\ref{rec:pointwise-smoothing}. Since
$\mathbb M(g)>0$, all these quantities are positive, and
\[
 H_\gamma(0)=0,
 \qquad H_\gamma(1)=\log Y_\gamma(g_\gamma).
\]
At a revelation in this line, the log ratio of its conditional
prediction to its uniform prediction is the corresponding
increment of $H_\gamma$. The factorization of $\pi_x$ shows
that $\log(\pi_x/\pi_{0,x})$ at the actual outcome is the sum
of the row and column increments.

Put $u_t=1-p_t$. From
\eqref{rec:smoothed-epsilon-schedule}, the weighted sum of the
increments on line $\gamma$ is exactly
\begin{align}
 C_\gamma(g,\tau)
 &:=\sum_{v\in\gamma}\epsilon(1-\tau_v)\,
                      \Delta H_\gamma(\tau_v)\notag\\
 &=\frac8{\sqrt\ell}\log Y_\gamma(g_\gamma)
   +\frac8\ell\sum_{t=0}^T
       \log\frac{Y_\gamma(g_\gamma)}
        {(Q_{A_\gamma(u_t)}Y_\gamma)(g_\gamma)}.
       \label{rec:smoothed-telescoping}
\end{align}
For each threshold, the increments occurring after $u_t$ sum
to $H_\gamma(1)-H_\gamma(u_t)$, which proves this identity.
The schedule is kept constant after its final jump all the way
to the last revelation. In particular, discarding the small
interval $p<p_{\min}$ in \eqref{rec:smoothed-slot-saving}
has not discarded any term in the telescoping identity.

Conditional on the full array $g$, each
$A_\gamma(u_t)$ is an independent-site Bernoulli-$p_t$ subset
of the domain of this line. The subsets for different $t$ are
nested, but their independence from $g$ and their individual
Bernoulli laws are all that is needed. The pointwise estimate
\eqref{rec:smoothing-ratio} and
\eqref{rec:smoothed-line-densities} yield
\begin{equation}\label{rec:smoothed-line-charge}
 \mathbb E_\tau C_\gamma(g,\tau)
 \le\frac{16}{\sqrt\ell}\log d_\gamma
          +C+C\ell^{-4}j_\gamma\log q_\gamma.
\end{equation}
Here $(8/\ell)(T+1)$ is bounded by an absolute constant.
For a trivial line $Y_\gamma=1$, the exact charge is zero,
so no constant is charged to that line. This is the point at
which smoothing before conditioning on success is essential:
\eqref{rec:smoothing-ratio} holds for the fixed $g_\gamma$,
although $g$ is now distributed according to $\mathbb N$.

Take expectations in \eqref{rec:smoothed-mixed-KL}, sum over
the cells, and use \eqref{rec:smoothed-exposure-KL}.
Conditional expectation replaces the $q_x$ expectation by the
actual revealed value. Its total weighted log ratio is
$\sum_\gamma C_\gamma(g,\tau)$. Equations
\eqref{rec:smoothed-slot-saving} and
\eqref{rec:smoothed-line-charge} therefore give
\begin{equation}\label{rec:smoothed-success-before-absorption}
 \log z\le-n+Cn m^{-1/40}
  +\sum_{\gamma:\,\mu_\gamma\ne(q_\gamma)}
       \left(\frac{16}{\sqrt\ell}\log d_\gamma
                  +C+C\ell^{-4}j_\gamma\log q_\gamma\right).
\end{equation}
All conditioning and normalization terms are retained in this
inequality; in particular there is no assumption that the line
laws remain independent after conditioning on $\mathcal S$.

For completeness, the dimension estimate used to absorb the last
two terms is
\begin{equation}\label{rec:smoothed-dimension-absorption}
 \log D_\mu\ge c j,
 \qquad \log D_\mu\ge c\log q,
 \quad j=q-\mu_1\ge1,\quad \ell(\mu)\le q/2,
\end{equation}
for sufficiently large $q$.
The first bound is \eqref{rec:band-dimension-lower}, applied at
size $q$. If $1\le j\le q/4$, \eqref{eq:near-row-hooks} gives
$D_\mu\ge e^{-1/2}\binom qj\ge e^{-1/2}q$, proving the second.
If $j>q/4$, the first bound gives $\log D_\mu\ge cq/4$ and
hence the second bound as well, after changing the absolute constant.

Since $q_\gamma\in[m,2m]$, that estimate implies
\[
 1+\ell^{-4}j_\gamma\log q_\gamma
            \le \frac{C}{\sqrt\ell}\log d_\gamma
               \qquad(\mu_\gamma\ne(q_\gamma)).
\]
It absorbs both the constants in
\eqref{rec:smoothed-success-before-absorption} and the cost
$N_*\log2$ in \eqref{rec:smoothed-majorization-cost}.
We conclude that
\begin{equation}\label{rec:smoothed-success-final}
 \log\mathbb P_\eta(\mathcal S)
 \le-n+Cn m^{-1/40}
                 +\frac{C}{\sqrt\ell}
                        \sum_{\gamma\in\mathcal L}\log d_\gamma.
\end{equation}
Finally $|K|=(b!)^a$, $|L|=(a!)^b$, and Stirling's estimate gives
\[
 \log\frac{n!}{(b!)^a(a!)^b}
 \le n+C(a+b)\log n
 \le n+Cn m^{-1/40}
\]
for sufficiently large $m$. Substitution in
\eqref{rec:smoothed-trace-probability} cancels the $-n$ in
\eqref{rec:smoothed-success-final}. Since
$\sum_\gamma\log d_\gamma\le D_0$, the result is
\[
 \operatorname{Tr}_{\rm reg}|JI|^4
 \le\exp\left\{\left(1+\frac{C}{\sqrt{\log m}}\right)D_0
                         +Cn m^{-1/40}\right\}.
\]
The local ranks have already been included exactly through
$\mathbb E|f_\gamma|^2=d_\gamma r_\gamma$ in
\eqref{rec:smoothed-trace-probability}; there is no additional
rank or regular-multiplicity factor to insert. For bounded
$m\ge2$, the same conclusion follows by enlarging $C$, using
$\operatorname{Tr}_{\rm reg}|JI|^4\le n!$. The case $m=1$
belongs to the finite base cases of the recursion rather than
to an estimate with denominator $\sqrt{\log m}$.
\end{proof}

\subsection{A sparse estimate from chronological mergers}

The entropy estimate will handle large-dimensional parent blocks.
For small diagram levels we use a separate sparse estimate, obtained
from a chronological merger forest and a bound on exceptional columns.
For each fixed $\delta>0$, there are $c_\delta>0$ and
$n_\delta<\infty$ such that, for every dyadic $n\ge n_\delta$
and every $\lambda\vdash n$,
\begin{equation}\label{rec:forest-sparse-general}
 0\le k=n-\lambda_1\le n^{1-\delta}
 \quad\Longrightarrow\quad
 \|T_n|_{V_\lambda}\|_{\op}\le n^{-c_\delta k}.
\end{equation}
\begin{proof}
The case $k=0$ again gives $1\le1$.
Use the tuple space, partial kernels $Q_I$, and centered
kernel $Z$ from
\eqref{rec:sparse-partial-kernel}--\eqref{rec:sparse-cover-majorant}.
The proof below estimates $Z$ directly by good rows and bad
columns; it does not use the preceding two-sweep chunk bound.
The case $k=1$ again has $Z=0$.
For $\delta\ge1$ this is the only possible positive level
for sufficiently large $n$, so assume $0<\delta<1$ and $k\ge2$.
Choose a fixed $0<\zeta<\delta/8$.

Fix an injective starting tuple $x$ and a subset $I$.
Use the unreduced probability experiment for the row of
$Q_I$: a common sweep for $I$ and independent fair-bit
paths for its complement, without conditioning on distinct
final sites.  Coordinates are numbered in their chronological
order.  For $u\ne v$, define
\[
 w_{uv}(x)=
 \sum_{\substack{1\le r\le d:\\
 (x_u)_{r+1:d}=(x_v)_{r+1:d}}}2^{-(r-1)}.
\]
Sharing the switch at time $r$ requires the displayed suffix
agreement and agreement of the two already updated prefixes
of length $r-1$.  Scan all such contacts in chronological
order, breaking ties by a fixed order of label pairs.
Retain a contact precisely when its endpoints are in different
components of the contacts already scanned.  The retained
edges form a forest.  On $\Gamma$ this forest has no isolated
vertices.

We need an unnormalized comparison for this exploration.
For any fixed forest with prescribed merger times $r_e$
satisfying the requisite initial suffix agreements,
\begin{equation}\label{rec:forest-history-mass}
 \Pr\{\text{the retained merger history is this forest}\}
 \le\prod_e2^{-(r_e-1)}.
\end{equation}
Here is a construction proving it without a conditional
independence assertion after nonmeeting events.
Start with one component for every label.  In a reference
experiment different components use independent switch
fields for their $I$-labels, and independent fair bits for
their other labels.  At each prescribed merger, require
only equality of the two indicated updated prefixes and
join the component measures.  Discard any configuration in
which two $I$-labels in the joined component occupy the same
site.  Then evolve the joined component with a common switch
field for its $I$-labels and independent bits for its remaining
labels.  Several prescribed mergers at the same time are
performed in the fixed tie order before updating that bit.
Do not impose the conditions that the distinct components
failed to meet at earlier times.  The reference measures
are subprobability measures, with their lost mass retained
in the calculation.

Just before time $r$, each component measure is invariant
under adding the same vector of $\mathbb F_2^{r-1}$ to all
its updated prefixes.  This follows inductively.
Initially there are no updated bits.  A required equality
of two prefixes preserves invariance under a common shift
of the joined component, as does the restriction excluding
coincident $I$-positions.  A switch update preserves
equivariance under shifts of the earlier coordinates and
is invariant under flipping the new output bit of every
label in the component.  The latter operation flips fair
coins and preserves the requirement that two jointly moved
labels leaving one switch use opposite outputs.
It supplies invariance in the newly updated coordinate.
Thus the invariant holds through all updates and mergers.

Before a prescribed merger the two component measures are
independent as unnormalized measures in the reference
construction.  If they have masses $a_1,a_2$, invariance
makes the indicated endpoint prefix uniform in each,
with total mass $a_i$.  The mass after requiring equality
is exactly $a_1a_2\,2^{-(r-1)}$.
Discarding invalid $I$-configurations can only decrease it.
All intervening Markov updates preserve mass.
Multiplying over mergers proves that the reference
subprobability has mass at most the right side of
\eqref{rec:forest-history-mass}.

To compare it with the original law, restrict to histories
having exactly the specified retained forest.  Until a
prescribed merger no switch is shared by labels in distinct
current components; otherwise the chronological scan would
have merged those components.  After all mergers at a given
time, every shared switch has its labels in one component.
On these histories the original and reference transition
weights therefore agree: a used switch of the joint
$I$-motion contributes one fair coin in either construction,
and every other label contributes its own fair bit.
The original history also satisfies all the prefix and
injectivity restrictions of the reference construction.
Its probability is thus bounded by the total reference
mass.  This proves \eqref{rec:forest-history-mass}.
In particular we have never divided by the probability of
having avoided earlier meetings; such a division would not
justify a uniform-prefix assertion.

Summing the time choices in
\eqref{rec:forest-history-mass} gives
\begin{equation}\label{rec:forest-cover-sum}
 \Pr_{Q_I(x,\cdot)}(\Gamma)
 \le\sum_{\substack{F\text{ a forest on }[k]\\
                         F\text{ has no isolated vertex}}}
                     \prod_{\{u,v\}\in F}w_{uv}(x).
\end{equation}
Keeping the final injectivity restriction would only
decrease the left side.  The estimate is uniform in $I$.

A suffix block of size $b=2^r$ fixes coordinates $r+1,\ldots,d$.
Call it dense for $x$ if its occupancy is at least two and
at least $b n^{-\zeta}$.  Call a vertex bad when it lies
in some dense block.  A row is good when fewer than $k/8$
vertices are bad.  Since the starting sites are distinct,
for every vertex $u$ one has
\[
 \sum_{v\ne u}w_{uv}(x)\le
       \sum_{r=1}^d2^{-(r-1)}2^r=2d.
\]
If $u$ is not bad, every block in this sum containing a
partner has occupancy less than $2^r n^{-\zeta}$, so
\begin{equation}\label{rec:forest-good-row-sum}
 \sum_{v\ne u}w_{uv}(x)\le2d n^{-\zeta}.
\end{equation}
Root each tree of a forest in
\eqref{rec:forest-cover-sum} and orient its edges toward
the root.  There are at most $k/2$ roots, because no tree
is a singleton.  On a good row at least $3k/8$ good
vertices are therefore nonroots.  Choose the root set
in at most $2^k$ ways and sum the parent of each nonroot
independently, dropping acyclicity and the other
restrictions to get an upper bound.  The row sums above
give the explicit rooted-forest estimate
\[
 \Pr_{Q_I(x,\cdot)}(\Gamma)
 \le2^k(2d)^k n^{-3\zeta k/8}
 \le n^{-\zeta k/4}
\]
for sufficiently large $n$.  This bound uses the merger
comparison, rather than independence of all contact events.

We next bound every column on the set $\mathcal B$ of bad
rows.  The transpose of $Q_I$ is a reversed joint sweep on
$|I|$ labels together with $k-|I|$ independent reversed
walkers, restricted to injective outputs.
This is exactly the partially coupled experiment in
Lemma~\ref{grid:occupation}.  In particular, for each set
of distinct sites $A$, its unrestricted endpoint law satisfies
\begin{equation}\label{rec:forest-mixed-occupation}
 \Pr\{\text{every site of }A\text{ is occupied}\}
 \le(k/n)^{|A|}.
\end{equation}
Use singleton test sets and $h_i=1$ in that lemma: the
product of endpoint counts is at least one on the event.
Thus the comparison applies also when independent walkers
can coincide.  Imposing distinct outputs can only lower
the column mass we are estimating.

The maximal dense suffix blocks of an injective bad row
are disjoint, because suffix blocks form a laminar family.
They contain at least $k/8$ occupied sites in total.
Put $M=n^\zeta$ and $q=Mk/n$.
For a chosen disjoint family, specify its occupied subsets
$A_B$, with $j_B=|A_B|\ge2$ and $j_B\ge |B|/M$.
By \eqref{rec:forest-mixed-occupation}, the probability of
occupying their union is at most $(k/n)^{\sum_Bj_B}$.
The condition $\sum_Bj_B\ge k/8$ allows multiplication
by $M^{\sum_Bj_B-k/8}\ge1$.
Summing the specified subsets, then dropping disjointness
and maximality of the block family, gives
\begin{equation}\label{rec:forest-bad-column-generating}
 \sum_{x\in\mathcal B}Q_I(x,y)
 \le M^{-k/8}\exp\left\{
   \sum_{\substack{b\text{ dyadic}\\1\le b\le n}}
     \frac nb
     \sum_{\substack{2\le j\le b\\j\ge b/M}}
                   \binom bj q^j\right\}.
\end{equation}
Indeed the sum over all families is bounded by the
product, over blocks, of one plus their total subset
weight, and this product is at most the displayed
exponential.  No independence of different occupied
blocks is being asserted.

We bound that exponent explicitly.  Put
$\varepsilon=eMq=eM^2k/n=o(1)$.
For $b\le2M$, $bq=o(1)$, and
$\sum_{j\ge2}\binom bj q^j\le C(bq)^2$.
Summing these dyadic sizes contributes at most
$Cnq^2M$.  For $b>2M$, let $j_0=\lceil b/M\rceil$.
The successive terms $(bq)^j/j!$ have ratio at most
$Mq$ for $j\ge j_0$.  Hence their tail is at most
$2(eMq)^{j_0}=2\varepsilon^{j_0}$ for large $n$.
The dyadic values of $b/M$ in this range at least double,
so
\[
 \sum_{b>2M}\frac nb\,2\varepsilon^{\lceil b/M\rceil}
 \le C\frac nM\varepsilon^2
 \le CnMq^2.
\]
Thus the whole exponent in
\eqref{rec:forest-bad-column-generating} is at most
\[
 CnMq^2=C M^3 k^2/n
 \le Ck n^{3\zeta-\delta}=o(k).
\]
Since $\zeta<\delta/8$, the last assertion is uniform
over the required range of $k$.  It follows that
\begin{equation}\label{rec:forest-small-bad-column}
 \sup_y\sum_{x\in\mathcal B}Q_I(x,y)
 \le n^{-\zeta k/16}
\end{equation}
for all sufficiently large $n$, uniformly in $I$.

Finally split $Z=Z_{\mathrm g}+Z_{\mathrm b}$ according
to whether its row belongs to $\mathcal B$.
Every $Q_I$ has row and column sums at most one.
The good-row coverage estimate and
\eqref{rec:sparse-cover-majorant} therefore give absolute
row sums at most $2^k n^{-\zeta k/4}$ and column sums at
most $2^k$ for $Z_{\mathrm g}$.
Equation \eqref{rec:forest-small-bad-column} gives column
sums at most $2^k n^{-\zeta k/16}$ and row sums at most
$2^k$ for $Z_{\mathrm b}$.
Schur's row-column bound yields
\[
 \|Z\|_{\op}
 \le2^k n^{-\zeta k/8}+2^k n^{-\zeta k/32}
 \le n^{-\zeta k/64}
\]
after increasing the starting size.
There is no additional falling-factorial normalization:
the $Q_I$ throughout are probability-mass kernels on
$\mathcal X_k$, and the uniform tuple inner product
differs from counting measure by a single constant.
Restriction to the $\lambda$-part in
\eqref{rec:sparse-centered-kernel} proves
\eqref{rec:forest-sparse-general}, with
$c_\delta=\zeta/64$.
\end{proof}

\subsection{Closing the singular-value recursion}

We combine the fourth-overlap estimate with the sparse norm bound.
The child exponents will enter only through positive powers; all
large-size thresholds are chosen independently of their values.

\begin{theorem}\label{rec:smoothed-singular}
There is a bounded sequence $p_d\ge8$ such that
$s_r(T_n|_{V_\lambda})\le(D_\lambda r)^{-1/p_d}$ for every
$n=2^d$, every $\lambda\vdash n$, and $1\le r\le D_\lambda$.
One may take, above a fixed base dimension,
\[
 p_d=(1+A/\sqrt d)\max\{p_{\lfloor d/2\rfloor},
                              p_{\lceil d/2\rceil}\}.
\]
\end{theorem}
\begin{proof}
We use rank-index bands, so that the number of bands and their
normalizations do not depend on the exponent inherited from the
children. For a child irreducible of dimension $D_\mu$, partition
its positive singular values into the index intervals
$[2^b,2^{b+1}-1]$. A nonempty band has rank $r\le2^b$ and norm at
most $(D_\mu2^b)^{-1/p'}\le(D_\mu r)^{-1/p'}$, where
$p'=\max(p_{\lfloor d/2\rfloor},p_{\lceil d/2\rceil})$.
The zero singular subspace contributes nothing. The row bands give
range projections $I$ for $X$, and the column bands give domain
projections $J$ for $Y$, in the factorization $T_n=YX$.

A profile chooses one type and one band on every row and column.
There are at most $\exp(C\sqrt q)$ types at a line of size $q$
and at most $1+\log_2(q!)\le Cq\log(q+1)$ index bands per type.
Thus the number $N_d$ of profiles satisfies
\begin{equation}\label{rec:smoothed-profile-count}
 \log N_d\le C\{a\sqrt b+b\sqrt a+(a+b)\log n\}
 \le Cn^{3/4}\log n.
\end{equation}
Types with more than half as many rows as line sites have zero
sweep operator by matching annihilation, and are omitted. For a
remaining profile put $D_0=\sum_i\log(D_{\mu_i}r_i)$.

Fix a nontrivial surviving parent type $\lambda$ and write
$L=\log D_\lambda$. Lemma~\ref{rec:smoothed-grid-fourth}, together
with regular multiplicity, gives
\[
 a_\lambda:=\operatorname{Tr}|(JI)_\lambda|^4
 \le A_0/D_\lambda,\qquad
 A_0=\exp\{(1+\eta)D_0+E_d\},\quad
 \eta=C_1/\sqrt d,
\]
where $E_d=n^{1-1/200}$ absorbs $Cn m^{-1/40}$ for all sufficiently
large $d$. The compression $(JI)_\lambda$ is a contraction. If its
singular values are $u_j$, then
$\sum_j u_j^8\le\sum_j u_j^4=a_\lambda$ and
$\sum_j u_j^8\le(\sum_j u_j^4)^2=a_\lambda^2$. Consequently
\begin{equation}\label{rec:smoothed-eighth}
 D_\lambda\operatorname{Tr}|(JI)_\lambda|^8
 \le A_0\min\{1,A_0/D_\lambda\}.
\end{equation}

Set $p_d=(1+A/\sqrt d)p'$ and $h=p_d/8\ge1$. Resolve
$|Y|^h|X^*|^h$ into its profile terms. A term has the form
$B_J(JI)A_I$, with $A_I,B_J$ supported on the indicated projections.
The product of their operator norms is at most $e^{-hD_0/p'}$.
The ideal property of Schatten norms and
\eqref{rec:smoothed-eighth} therefore show that the logarithm of
its eighth power norm, after multiplication by $D_\lambda$, is
at most
\begin{equation}\label{rec:smoothed-profile-saving}
 -\frac{A-C_1}{\sqrt d}D_0+E_d
 +\min\{0,(1+\eta)D_0+E_d-L\}.
\end{equation}
This bound is independent of the size of $p'$.

Choose $A>C_1+8$. In the dense range
$k>n^{1-1/400}$, \eqref{rec:band-dimension-lower} gives $L\ge c n^{1-1/400}$. Hence $E_d=o(L/\sqrt d)$ and
$\log N_d=o(L/\sqrt d)$, with thresholds independent of $p'$.
If $D_0\ge L/4$, the first two terms in
\eqref{rec:smoothed-profile-saving} are at most $-c_1L/\sqrt d$.
If $D_0<L/4$, then $(1+\eta)D_0+E_d\le L/2$ for large $d$,
and the same expression is at most $-L/3$.
Thus every profile term $F_\gamma$ satisfies
\[
 D_\lambda\|F_\gamma\|_8^8
 \le e^{-c_2L/\sqrt d}.
\]
The triangle inequality, followed by
\eqref{rec:smoothed-profile-count}, gives
\begin{equation}\label{rec:smoothed-powered-moment}
 D_\lambda\big\||Y|^h|X^*|^h\big\|_8^8
 \le N_d^8e^{-c_2L/\sqrt d}\le1.
\end{equation}
All norms in this paragraph are on the single parent irreducible.

Apply Lemma~\ref{lem:positive-power-comparison} with $q=8$:
\begin{equation}\label{rec:smoothed-power-interpolation}
 \|YX\|_p\le
 \big\||Y|^{p/8}|X^*|^{p/8}\big\|_8^{8/p},\qquad p\ge8.
\end{equation}
Combining this with \eqref{rec:smoothed-powered-moment} yields
$\sum_r s_r(T_n|_{V_\lambda})^{p_d}\le D_\lambda^{-1}$.
Since the singular values decrease, their first $r$ terms imply
the required bound $(D_\lambda r)^{-1/p_d}$.

In the sparse range $1\le k\le n^{1-1/400}$, use
\eqref{rec:forest-sparse-general} with $\delta=1/400$ and
$\log(D_\lambda r)\le2k\log n$. It suffices to keep all exponents
at least $2/c_{1/400}$, which is achieved by increasing the common
finite base exponent. Lemma~\ref{lem:finitegap} supplies a
sufficiently large common base exponent for the finitely many
nonconstant blocks and singular indices at those sizes. The trivial block has $D_\lambda=r=1$
and singular value one, and the killed blocks have singular value
zero, so both satisfy the statement directly.

Finally the large-size threshold depends only on the absolute
constants in the overlap and dimension bounds, not on the finite base
exponent. Lemma~\ref{lem:dyadic-exponents}, with $\gamma=1/2$,
bounds the product of $1+A/\sqrt d$ along every rounded-halving
branch. Hence $\sup_d p_d<\infty$.
\end{proof}

The indexed power bounds of Sections~\ref{sec:convex-grid}--\ref{sec:smoothed-grid}
now give the same fixed-sweep mixing consequence by the conversion in
Section~\ref{sec:spectral-conversion}. Their usable operator exponents
are, respectively, $10^{-8}c$, $\alpha$, and $1/\sup_dp_d$.
For any of these exponents $a>0$, an integer $R$ with $2aR\ge3$
makes the regular contribution
$D_\lambda\|T_n(\lambda)^R\|_{\HS}^2$ of each nontrivial
nonzero block at most $D_\lambda^{-1}$. The sign block vanishes,
and Lemma~\ref{lem:degrees} makes the sum tend to zero. Thus $R$
forward sweeps, or $Rd$ physical shuffles, converge to uniform in
total variation. The matching order lower bound is
Lemma~\ref{lem:support}.
